\subsection{Raumverwaltung und Belegungsdaten}

Die Daten der Raumverwaltung werden nicht manuell gepflegt, sondern mit einem Python-Skript aus
bereits vorhandenen Stammdaten abgeleitet. Dadurch bleiben Stationen, Zimmer, Pflegepersonal und
Buchungen konsistent zueinander. Das Skript erzeugt bzw. aktualisiert die CSV-Dateien für
Stationen, Räume, Pflegekräfte und Buchungen. Bei dem zuletzt erzeugten Datensatz wurden
18 Stationen, 184 Zimmer und 8918 Buchungseinträge generiert.

\subsubsection*{Eingabedaten und Abhängigkeiten}

Als Grundlage verwendet das Skript vier CSV-Dateien:

\begin{itemize}
    \item \texttt{departments.csv}: enthält die medizinischen Abteilungen sowie das Gebäude, aus dem später die Etage der Räume abgeleitet wird.
    \item \texttt{employees.csv}: ordnet Mitarbeitende ihren Abteilungen zu.
    \item \texttt{nurses.csv}: liefert die Menge der Pflegekräfte, die anschließend einer Station zugeordnet werden.
    \item \texttt{patients.csv}: enthält die Patienten-IDs, die für die späteren Buchungen benötigt werden.
\end{itemize}

Wichtig ist, dass diese Dateien vor der Ausführung bereits vollständig und zueinander passend
vorliegen. Besonders relevant sind dabei die Fremdschlüsselbeziehungen aus dem Datenbankschema:
\texttt{station.department} verweist auf \texttt{department.id}, \texttt{rooms.station} verweist
auf \texttt{station.id}, \texttt{nurses.station} verweist auf \texttt{station.id} und
\texttt{bookings.room} sowie \texttt{bookings.patient} verweisen auf \texttt{rooms.id} bzw.
\texttt{patient.id}. Deshalb müssen die Abteilungs-IDs, Mitarbeiter-IDs, Stations-IDs, Raum-IDs
und Patienten-IDs in den erzeugten CSV-Dateien exakt mit den referenzierten Tabellen übereinstimmen.

\subsubsection*{Erzeugung der Stationen und Räume}

Im Skript ist mit \texttt{INPATIENT\_PROFILES} festgelegt, welche Abteilungen stationäre
Einheiten erhalten. Nicht jede Abteilung wird automatisch zu einer Station, sondern nur die dort
aufgeführten Abteilungen. Für jede dieser Abteilungen wird ein Stationsprofil definiert. Dieses
Profil enthält den Stationsnamen, den Stationstyp und einen Faktor \texttt{beds\_per\_nurse}.
Der Stationstyp unterscheidet beispielsweise Intensivstation, Allgemeinstation, Chirurgie,
Innere Medizin, Onkologie, Pädiatrie und Psychiatrie.

Die Anzahl der Betten wird aus der Anzahl der Pflegekräfte einer Abteilung berechnet:
Die Mitarbeitenden werden zuerst über \texttt{employees.csv} ihrer Abteilung zugeordnet.
Anschließend zählt das Skript, wie viele Einträge aus \texttt{nurses.csv} zu jeder Abteilung
gehören. Aus dieser Anzahl und dem Profilfaktor entsteht die Bettenzahl der Station. Damit keine
Station unrealistisch klein wird, wird mindestens mit sechs Betten gerechnet.

Aus der Bettenzahl werden anschließend konkrete Zimmer erzeugt. Die Funktion
\texttt{build\_room\_sizes()} verteilt die Betten auf Ein- und Zweibettzimmer. Das Verhältnis ist
pro Stationstyp unterschiedlich festgelegt. Intensivstationen bestehen vollständig aus
Einbettzimmern, während andere Stationstypen überwiegend Zweibettzimmer enthalten und nur einen
definierten Anteil an Einzelzimmern haben. Dieses Verhältnis ist wichtig, weil es später bestimmt,
wie viele Patienten parallel einem Raum zugeordnet sein können.

Jeder Raum erhält eine technische \texttt{id}, eine \texttt{station}, eine sichtbare
\texttt{number}, eine \texttt{floor}-Angabe und die Anzahl der \texttt{beds}. Die Etage wird aus
dem ersten Zahlenzeichen des Gebäudefeldes der Abteilung abgeleitet, z.\,B. aus \texttt{4B} die
Etage \texttt{4}. Für Verknüpfungen ist jedoch nicht die Raumnummer entscheidend, sondern die
technische Raum-ID. Die Raumnummer dient vor allem der Anzeige, während \texttt{rooms.id} der
Primärschlüssel ist und in \texttt{bookings.room} referenziert wird.

\subsubsection*{Zuordnung der Pflegekräfte}

Neben den Stationen und Räumen schreibt das Skript auch \texttt{nurses.csv} neu. Dabei bleibt die
ID der Pflegekraft erhalten, die Spalte \texttt{station} wird jedoch aus der Abteilung des
zugehörigen Employee-Eintrags bestimmt. Dadurch entsteht eine direkte Beziehung zwischen
Pflegekraft und Station. Dieser Schritt ist notwendig, weil die Datenbank für Pflegekräfte einen
Fremdschlüssel auf \texttt{station.id} vorsieht. Pflegekräfte aus Abteilungen ohne stationäres
Profil erhalten keine Stationszuordnung.

\subsubsection*{Generierung der Buchungen}

Die Buchungen werden aus den erzeugten Zimmern und der Patientenliste aufgebaut. Zuerst werden alle
Patienten-IDs zufällig gemischt. Danach wird für jedes Bett eines Raums eine eigene Zeitachse
simuliert. Dadurch kann ein Zweibettzimmer zwei parallele Buchungen haben, aber nicht beliebig viele.
Die Aufenthaltsdauer wird abhängig vom Stationstyp aus einer Dreiecksverteilung bestimmt. Eine
psychiatrische Station erzeugt dadurch längere Aufenthalte als beispielsweise eine Allgemeinstation
oder eine chirurgische Station.

Der betrachtete Zeitraum umfasst im Skript 180 Tage in der Vergangenheit und 14 Tage in der Zukunft.
Zwischen zwei Aufenthalten werden zusätzlich kleine Lücken eingefügt, die ebenfalls vom Stationstyp
abhängen. Für jeden erzeugten Aufenthalt wird ein Patient aus dem Patientenpool entnommen und nur
einmal verwendet. Dadurch werden doppelte oder zeitlich überlappende Buchungen desselben Patienten
vermieden. Falls nicht genug Patienten vorhanden sind, bricht das Skript mit einer Fehlermeldung ab.
Diese Prüfung ist wichtig, weil sonst Buchungen entstehen könnten, deren \texttt{patient}-Wert nicht
mehr sauber auf die Patiententabelle verweist.

Der Status einer Buchung wird anhand des Aufnahme- und Entlassdatums relativ zum Ausführungsdatum
des Skripts bestimmt. Aktuelle Aufenthalte erhalten \texttt{checked\_in}, zukünftige Aufenthalte
\texttt{confirmed} oder \texttt{pending}. Bereits vergangene Aufenthalte werden überwiegend als
\texttt{completed} markiert, ein kleiner Anteil als \texttt{checked\_out\_early}. Zusätzlich werden
separate Ereignisse für \texttt{cancelled} und \texttt{no\_show} erzeugt. Diese Ereignisse zählen
nicht als reale Belegung, bilden aber typische Sonderfälle für die API und das Frontend ab.

\subsubsection*{Integritätsregeln beim Import}

Beim Import der CSV-Dateien muss die Reihenfolge der Tabellen beachtet werden. Zuerst müssen
\texttt{department}, \texttt{person}, \texttt{patient} und \texttt{employee} vorhanden sein. Danach
können \texttt{station} und \texttt{rooms} importiert werden. Erst anschließend sind
\texttt{nurses} und \texttt{bookings} sicher importierbar, weil ihre Fremdschlüssel dann bereits
existieren.

Außerdem müssen folgende Bedingungen eingehalten werden:

\begin{itemize}
    \item Jede \texttt{station.department}-ID muss in \texttt{department.id} existieren.
    \item Jede \texttt{rooms.station}-ID muss in \texttt{station.id} existieren.
    \item Jede \texttt{nurses.station}-ID muss entweder leer sein oder auf eine vorhandene Station zeigen.
    \item Jede \texttt{bookings.room}-ID muss auf einen vorhandenen Raum zeigen.
    \item Jede \texttt{bookings.patient}-ID muss auf einen vorhandenen Patienten zeigen.
    \item Pro belegungsrelevantem Zeitraum darf die Anzahl paralleler Buchungen eines Raums die Bettenzahl des Raums nicht überschreiten.
    \item Die Werte in \texttt{bookings.state} müssen zum Booking-State-Enum der Datenbank passen.
\end{itemize}

Durch den festen Zufallsseed \texttt{RANDOM\_SEED = 42} bleiben die generierten Daten bei gleichen
Eingabedateien reproduzierbar. Das ist für Tests und Dokumentation hilfreich, weil die Raumstruktur
und die erzeugten Buchungen nicht bei jedem Durchlauf unkontrolliert variieren. Gleichzeitig bilden
die Zufallsverteilungen genug Varianz ab, um realistische Belegungszustände, freie Betten,
abgeschlossene Aufenthalte und Sonderfälle wie Stornierungen testen zu können.
